\begin{lemma}
For fixed $g$ and $k<g(k)$, the map $f\mapsto\sigma_f$ is
prefix-continuous.  Explicitly, for every $f$ and every output length $n$,
some finite prefix of $f$ determines all values
$\sigma_f(l)$ for $l<n$.
\end{lemma}
\begin{proof}
Fix $f$ and $n$, and write
\[
G(i)=\mathsf{OrbitPoint}(g,k,i).
\]
By \noderef{OrbitBlockCoverage}, for each $l<n$ either $l<G(0)$, or there
is a unique index $q_l$ such that
\[
G(q_l)\leq l<G(q_l+1).
\]
In the second case, the endpoint inequality makes
\[
E_l=\{i\in\mathbb N\mid l<G(i+1)\}
\]
nonempty.  Let $i_l$ be its least element.  The least-index operation in
the definition of \noderef{BlockSigma} selects exactly $i_l$.  Since the
predicate defining $E_l$ involves only $g,k,l$, this selected index depends
only on $g,k,l$, not on $f$; moreover, $i_l\leq q_l$ because $q_l\in E_l$.
The uniqueness in \noderef{OrbitBlockCoverage} makes the block witness
$q_l$ unambiguous for each $l$, so these are precisely the finitely many
block indices that must be bounded.

We now give an explicit finite bound for all indices that can be relevant
below the output bound $n$.  Define $B_0=0$ and process the integers
$r=0,\ldots,n-1$ recursively.  If $r<G(0)$, put $B_{r+1}=B_r$; otherwise
put
\[
B_{r+1}=\max\{B_r,q_r+1\}.
\]
The first clause records the initial interval and introduces no block index.
The second clause gives $q_r<B_{r+1}$.  Hence every block index $q_l$ for
$l<n$ satisfies $q_l<B_n$.  Set $m=B_n$ (with $m=0$ when $n=0$).

Let $h$ satisfy \(\mathsf{PrefixAgree}(m,f,h)\).  By the definition in
\noderef{PrefixAgree}, this says $f(i)=h(i)$ for every $i<m$.  Fix $l<n$.
If $l<G(0)$, both applications of \noderef{BlockSigma} take its initial
branch, so
\[
\mathsf{BlockSigma}(g,k,h,l)=l=\mathsf{BlockSigma}(g,k,f,l).
\]
Otherwise, the index selected in both applications of \noderef{BlockSigma}
is the same $i_l$, because its defining least-index predicate involves only
$g,k,l$.  Since $i_l\leq q_l<m$, \noderef{PrefixAgree} gives
$f(i_l)=h(i_l)$, and therefore
\[
h(i_l)-i_l=f(i_l)-i_l.
\]
The block branch of \noderef{BlockSigma} now gives
\[
\begin{aligned}
\mathsf{BlockSigma}(g,k,h,l)
  &=g^{[h(i_l)-i_l]}(l)\\
  &=g^{[f(i_l)-i_l]}(l)\\
  &=\mathsf{BlockSigma}(g,k,f,l).
\end{aligned}
\]
Thus the equality holds for every $l<n$, and the displayed $m$ proves the
claimed prefix-continuity statement.
\end{proof}
